\section{Trace 分析：数据驱动的改进循环}

\subsection{Trace Analyzer Skill}

LangChain 将 Trace 分析封装为一个可复用的 Agent Skill（Trace Analyzer），用于系统性地分析实验运行中的错误并生成改进建议。其工作流程如下：

\begin{enumerate}[leftmargin=2em]
    \item \textbf{获取 Traces}：从 LangSmith 拉取实验运行的执行追踪数据
    \item \textbf{并行错误分析}：派出多个子 Agent 并行分析不同 Trace 中的错误
    \item \textbf{汇总综合}：主 Agent 综合所有子 Agent 的发现，生成改进建议
    \item \textbf{针对性修改}：人工审核（可选）后，对 Harness 做针对性调整
\end{enumerate}

\begin{knowledgebox}{与 Boosting 的类比}
这种方法类似于机器学习中的 Boosting：每轮训练（Harness 迭代）重点关注上一轮中模型犯错的样本（失败的 Trace）。通过聚焦错误，迭代改进能快速收敛到更好的性能。
\end{knowledgebox}

\subsection{发现的常见失败模式}

通过 Trace 分析，LangChain 发现了以下 Agent 的常见失败模式：

\begin{table}[H]
\centering
\caption{Agent 常见失败模式及应对策略}
\begin{tabular}{p{4cm}p{4cm}p{4.5cm}}
\toprule
\textbf{失败模式} & \textbf{表现} & \textbf{应对策略} \\
\midrule
不自我验证 & 写完代码后只看一遍就提交 & Build-Verify 循环 \\
推理错误 & 逻辑链条出错 & 提升 Reasoning 预算 \\
忽视任务指令 & 未严格遵循 Task Spec & 强化 Prompt \\
超时 & 在次要问题上花费过多时间 & 时间预算管理 \\
死循环 & 对同一文件反复小修改 & LoopDetectionMiddleware \\
环境认知不足 & 不了解可用工具和目录结构 & 上下文注入 \\
\bottomrule
\end{tabular}
\end{table}

\begin{warningbox}{过拟合风险}
在 Trace 驱动的改进中，需要警惕\textbf{过拟合}：针对某个特定任务做出的修改可能导致其他任务性能回退。人工审核步骤（Step 3）的价值在于判断改进是否具有泛化性，而非仅仅修复个案。
\end{warningbox}

\subsection{本章小结}

Trace 分析是 Harness Engineering 的"眼睛"。通过将 Trace 分析自动化并封装为 Skill，LangChain 实现了快速、可重复的改进循环。发现的六类典型失败模式（不验证、推理错误、忽视指令、超时、死循环、环境认知不足）构成了后续所有 Harness 改进的方向指引。


